% GENERATED FILE. Edit canonical lessons, then run npm run generate:books.
% canonical-source-hash: fnv1a64:e36e05e7563a882b
% canonical-lessons: KA-C510-noyu, KA-C510-jomu, KA-C510-ubbu, KA-C510-osadu, KA-C510-davade

\chapter{Aches and Swellings}
\label{ch:ka-aches-and-swellings}
\hlchaptermodality{510}

\begin{chapteropening}
\textbf{By the end of this chapter:} \emph{I can say to ache, numbness, to swell, the gums and the jaw.}

Complete the last lesson of chapter 510: I can say to ache, numbness, to swell, the gums and the jaw.
\end{chapteropening}

\section[\texorpdfstring{\kn{ನೋಯು} (nōyu)}{ನೋಯು (nōyu)}]{\kn{ನೋಯು} (nōyu) — to ache; to be hurt}

\label{lesson:KA-C510-noyu}



\begin{quote}
Before the new one: say the Kannada for an arm, then the Kannada for to limp.
\end{quote}

\subsection*{You'll want to know: \kn{ನೋಯು}}
\textbf{\kn{ನೋಯು}} — \emph{nōyu} — \textquotedblleft{}to ache; to be hurt\textquotedblright{}.

\begin{quote}
\kn{ತುಂಬಾ} \kn{ನಡೆದಿದ್ದರಿಂದ} \kn{ನನ್ನ} \kn{ಕಾಲು} \kn{ನೋಯುತ್ತಿದೆ}.
\end{quote}

\begin{quote}
\emph{tumbā naḍediddarinda nanna kālu nōyuttide.}
\end{quote}

\begin{quote}
My leg aches because I walked a lot.
\end{quote}

Past (I …): \kn{ನೊಂದೆ}, \emph{nonde}.

The noun is \kn{ನೋವು}, pain. Feelings ache too: \kn{ಅವನ} \kn{ಮಾತಿನಿಂದ} \kn{ನಾನು} \kn{ನೊಂದೆ}, his words hurt me.

\begin{cousinweb}
\textbf{\kn{ನೋಯು}} (\emph{nōyu}) is native Dravidian: DEDR 3793 lists Kannada \emph{nōyu}, \textquotedblleft{}to feel pain, be grieved\textquotedblright{}, beside Tamil \emph{nō}. The noun \kn{ನೋವು}, pain, is built on it.
\end{cousinweb}

\begin{grammarlens}[title={What you've built}]
One more word for describing what you lived and felt.
\end{grammarlens}

\subsection*{Your turn}
\begin{itemize}
\raggedright
  \item \emph{Say it:} \emph{nōyu}
  \item \emph{Say it:} \emph{nōyu}, once more
  \item \emph{Say it:} \emph{tumbā naḍediddarinda nanna kālu nōyuttide.} — My leg aches because I walked a lot.
  \item \emph{Recall:} say the Kannada for an arm, then the Kannada for to limp, then say \emph{nōyu} again
\end{itemize}

\begin{tcolorbox}[breakable,colback=teal!4,colframe=teal!35!black,title={Before you move on}]
What does \kn{ನೋಯು} mean? (\textquotedblleft{}To ache; to be hurt\textquotedblright{}.) Say it once more.
\end{tcolorbox}
\section[\texorpdfstring{\kn{ಜೋಮು} (jōmu)}{ಜೋಮು (jōmu)}]{\kn{ಜೋಮು} (jōmu) — numbness; pins and needles}

\label{lesson:KA-C510-jomu}



\begin{quote}
Before the new one: say the Kannada for to limp, then the Kannada for to ache.
\end{quote}

\subsection*{You'll want to know: \kn{ಜೋಮು}}
\textbf{\kn{ಜೋಮು}} — \emph{jōmu} — \textquotedblleft{}numbness; pins and needles\textquotedblright{}.

\begin{quote}
\kn{ತುಂಬಾ} \kn{ಹೊತ್ತು} \kn{ಕುಳಿತಿದ್ದರಿಂದ} \kn{ನನ್ನ} \kn{ಕಾಲಿಗೆ} \kn{ಜೋಮು} \kn{ಹಿಡಿಯಿತು}.
\end{quote}

\begin{quote}
\emph{tumbā hottu kuḷitiddarinda nanna kālige jōmu hiḍiyitu.}
\end{quote}

\begin{quote}
From sitting too long, my leg went numb.
\end{quote}

\kn{ಜೋಮು} \kn{ಹಿಡಿ}, to go numb: the numbness \textquotedblleft{}grips\textquotedblright{} the limb.

\begin{cousinweb}
\textbf{\kn{ಜೋಮು}} (\emph{jōmu}) is native Dravidian and imitative: DEDR 2578 lists Kannada \emph{jum}, \emph{jummu} and also \emph{jommu}, \emph{jōmu} for the tingling numbness of a limb kept long in one position.
\end{cousinweb}

\begin{grammarlens}[title={What you've built}]
One more word for describing what you lived and felt.
\end{grammarlens}

\subsection*{Your turn}
\begin{itemize}
\raggedright
  \item \emph{Say it:} \emph{jōmu}
  \item \emph{Say it:} \emph{jōmu}, once more
  \item \emph{Say it:} \emph{tumbā hottu kuḷitiddarinda nanna kālige jōmu hiḍiyitu.} — From sitting too long, my leg went numb.
  \item \emph{Recall:} say the Kannada for to limp, then the Kannada for to ache, then say \emph{jōmu} again
\end{itemize}

\begin{tcolorbox}[breakable,colback=teal!4,colframe=teal!35!black,title={Before you move on}]
What does \kn{ಜೋಮು} mean? (\textquotedblleft{}Numbness; pins and needles\textquotedblright{}.) Say it once more.
\end{tcolorbox}
\section[\texorpdfstring{\kn{ಉಬ್ಬು} (ubbu)}{ಉಬ್ಬು (ubbu)}]{\kn{ಉಬ್ಬು} (ubbu) — to swell; to bulge}

\label{lesson:KA-C510-ubbu}



\begin{quote}
Before the new one: say the Kannada for to ache, then the Kannada for numbness.
\end{quote}

\subsection*{You'll want to know: \kn{ಉಬ್ಬು}}
\textbf{\kn{ಉಬ್ಬು}} — \emph{ubbu} — \textquotedblleft{}to swell; to bulge\textquotedblright{}.

\begin{quote}
\kn{ಜೇನುನೊಣ} \kn{ಕಚ್ಚಿದ} \kn{ಮೇಲೆ} \kn{ಅವನ} \kn{ಕೈ} \kn{ಉಬ್ಬಿತು}.
\end{quote}

\begin{quote}
\emph{jēnunoṇa kaccida mēle avana kai ubbitu.}
\end{quote}

\begin{quote}
After the bee stung him, his hand swelled.
\end{quote}

Past (it …): \kn{ಉಬ್ಬಿತು}, \emph{ubbitu}.

\kn{ಊತ} is a swelling; \kn{ಉಬ್ಬು} is the swelling up.

\begin{cousinweb}
\textbf{\kn{ಉಬ್ಬು}} (\emph{ubbu}) is native Dravidian: DEDR 666 lists Kannada \emph{ubbu}, \textquotedblleft{}to swell, rise up\textquotedblright{}, with the noun \emph{ubbuvike}, \textquotedblleft{}swelling\textquotedblright{}.
\end{cousinweb}

\begin{grammarlens}[title={What you've built}]
One more word for describing what you lived and felt.
\end{grammarlens}

\subsection*{Your turn}
\begin{itemize}
\raggedright
  \item \emph{Say it:} \emph{ubbu}
  \item \emph{Say it:} \emph{ubbu}, once more
  \item \emph{Say it:} \emph{jēnunoṇa kaccida mēle avana kai ubbitu.} — After the bee stung him, his hand swelled.
  \item \emph{Recall:} say the Kannada for to ache, then the Kannada for numbness, then say \emph{ubbu} again
\end{itemize}

\begin{tcolorbox}[breakable,colback=teal!4,colframe=teal!35!black,title={Before you move on}]
What does \kn{ಉಬ್ಬು} mean? (\textquotedblleft{}To swell; to bulge\textquotedblright{}.) Say it once more.
\end{tcolorbox}
\section[\texorpdfstring{\kn{ಒಸಡು} (osaḍu)}{ಒಸಡು (osaḍu)}]{\kn{ಒಸಡು} (osaḍu) — the gums}

\label{lesson:KA-C510-osadu}



\begin{quote}
Before the new one: say the Kannada for numbness, then the Kannada for to swell.
\end{quote}

\subsection*{You'll want to know: \kn{ಒಸಡು}}
\textbf{\kn{ಒಸಡು}} — \emph{osaḍu} — \textquotedblleft{}the gums\textquotedblright{}.

\begin{quote}
\kn{ಒಸಡಿನಿಂದ} \kn{ರಕ್ತ} \kn{ಬರುತ್ತಿದೆ} \kn{ಅಂತ} \kn{ಅವನು} \kn{ವೈದ್ಯರಿಗೆ} \kn{ಹೇಳಿದನು}.
\end{quote}

\begin{quote}
\emph{osaḍininda rakta baruttide anta avanu vaidyarige hēḷidanu.}
\end{quote}

\begin{quote}
He told the doctor that his gums were bleeding.
\end{quote}

Plural: \kn{ಒಸಡುಗಳು}, \emph{osaḍugaḷu}.

Often in the plural, \kn{ಒಸಡುಗಳು}.

\begin{cousinweb}
\textbf{\kn{ಒಸಡು}} (\emph{osaḍu}) is a native Kannada word for the gums; its origin is not settled.
\end{cousinweb}

\begin{grammarlens}[title={What you've built}]
One more word for describing what you lived and felt.
\end{grammarlens}

\subsection*{Your turn}
\begin{itemize}
\raggedright
  \item \emph{Say it:} \emph{osaḍu}
  \item \emph{Say it:} \emph{osaḍu}, once more
  \item \emph{Say it:} \emph{osaḍininda rakta baruttide anta avanu vaidyarige hēḷidanu.} — He told the doctor that his gums were bleeding.
  \item \emph{Recall:} say the Kannada for numbness, then the Kannada for to swell, then say \emph{osaḍu} again
\end{itemize}

\begin{tcolorbox}[breakable,colback=teal!4,colframe=teal!35!black,title={Before you move on}]
What does \kn{ಒಸಡು} mean? (\textquotedblleft{}The gums\textquotedblright{}.) Say it once more.
\end{tcolorbox}
\section[\texorpdfstring{\kn{ದವಡೆ} (davaḍe)}{ದವಡೆ (davaḍe)}]{\kn{ದವಡೆ} (davaḍe) — the jaw}

\label{lesson:KA-C510-davade}



\begin{quote}
Before the new one: say the Kannada for to swell, then the Kannada for the gums.
\end{quote}

\subsection*{You'll want to know: \kn{ದವಡೆ}}
\textbf{\kn{ದವಡೆ}} — \emph{davaḍe} — \textquotedblleft{}the jaw\textquotedblright{}.

\begin{quote}
\kn{ಆನೆಯ} \kn{ದವಡೆ} \kn{ತುಂಬಾ} \kn{ಬಲವಾಗಿದೆ}.
\end{quote}

\begin{quote}
\emph{āneya davaḍe tumbā balavāgide.}
\end{quote}

\begin{quote}
An elephant's jaw is very strong.
\end{quote}

\kn{ದವಡೆ} \kn{ಹಲ್ಲು}, a molar, a jaw tooth.

\begin{cousinweb}
\textbf{\kn{ದವಡೆ}} (\emph{davaḍe}) is a native Kannada word for the jaw; its origin is not settled.
\end{cousinweb}

\begin{grammarlens}[title={What you've built}]
One more word for describing what you lived and felt.
\end{grammarlens}

\subsection*{Your turn}
\begin{itemize}
\raggedright
  \item \emph{Say it:} \emph{davaḍe}
  \item \emph{Say it:} \emph{davaḍe}, once more
  \item \emph{Say it:} \emph{āneya davaḍe tumbā balavāgide.} — An elephant's jaw is very strong.
  \item \emph{Recall:} say the Kannada for to swell, then the Kannada for the gums, then say \emph{davaḍe} again
\end{itemize}

\begin{tcolorbox}[breakable,colback=teal!4,colframe=teal!35!black,title={Before you move on}]
What does \kn{ದವಡೆ} mean? (\textquotedblleft{}The jaw\textquotedblright{}.) Say it once more.
\end{tcolorbox}
